\documentclass{article}

% packages
\usepackage{amsmath, amsthm, thmtools, amsfonts, amssymb, luacode, catchfile, hyperref, ifthen, graphicx, bbm}
\ifcsname c@kobocompile\endcsname
	\usepackage[a5paper, total={1072pt, 1448pt}, margin=10pt, includeheadfoot]{geometry} % set page margins
\else
	\usepackage[a4paper, margin=50pt, includeheadfoot]{geometry}
\fi
\usepackage[shortlabels]{enumitem}
\usepackage[skip=3pt, indent=0pt]{parskip}

% language
\usepackage[bidi=basic, layout=tabular, provide=*]{babel}
\ifcsname c@english\endcsname
	\babelprovide[main, import]{english}
\else
	\babelprovide[main, import, alph=letters.gershayim]{hebrew}
	\setlist[enumerate,2]{label={\alph*.}}
	\babelprovide{rl}
\fi
\babelfont{rm}[Renderer=Harfbuzz]{Libertinus Serif}
\babelfont{sf}{Libertinus Sans}
\babelfont{tt}{Libertinus Mono}

% style
\ifcsname c@stoppagebreaks\endcsname
\else
	\AddToHook{cmd/section/before}{\clearpage}	% Add line break before section
\fi
\linespread{1.3}
\setcounter{secnumdepth}{0}		% Remove default number tags from sections, this won't do well with theorems
\AtBeginDocument{\setlength{\belowdisplayskip}{3pt}}
\AtBeginDocument{\setlength{\abovedisplayskip}{3pt}}
\graphicspath{ {../images/} }

% operators
\DeclareMathOperator\cis{cis}
\DeclareMathOperator\Sp{Span}
\DeclareMathOperator\tr{tr}
\DeclareMathOperator\im{Im}
\DeclareMathOperator\re{Re}
\DeclareMathOperator\diag{diag}
\DeclareMathOperator*\lowlim{\underline{lim}}
\DeclareMathOperator*\uplim{\overline{lim}}
\DeclareMathOperator\rng{rng}
\DeclareMathOperator\Sym{Sym}
\DeclareMathOperator\Arg{Arg}
\DeclareMathOperator\Log{Log}
\DeclareMathOperator\dom{dom}
\DeclareMathOperator\supp{Supp}
\DeclareMathOperator\var{Var}
\DeclareMathOperator\cov{Cov}
\DeclareMathOperator\End{End}
\DeclareMathOperator\id{id}
\DeclareMathOperator\cl{cl}
\DeclareMathOperator\gal{Gal}
\DeclareMathOperator\dist{dist}
\DeclareMathOperator\form{form}
\DeclareMathOperator\term{term}
\DeclareMathOperator\sent{sent}
\DeclareMathOperator\Th{Th}
\DeclareMathOperator\Hom{Hom}
\DeclareMathOperator\Aut{Aut}
\DeclareMathOperator\Char{char}
\DeclareMathOperator\vol{vol}
\DeclareMathOperator\GL{GL}
\DeclareMathOperator\SL{SL}
\DeclareMathOperator\lcm{lcm}

% commands
\newcommand{\cupdot}{\mathbin{\mathaccent\cdot\cup}}
\newcommand{\indicator}{\mathbbm{1}}
\newcommand{\AAA}[0]{\mathbb{A}}
\newcommand{\Aa}[0]{\mathcal{A}}
\newcommand{\BB}[0]{\mathbb{B}}
\newcommand{\Bb}[0]{\mathcal{B}}
\newcommand{\CC}[0]{\mathbb{C}}
\newcommand{\Cc}[0]{\mathcal{C}}
\newcommand{\DD}[0]{\mathbb{D}}
\newcommand{\Dd}[0]{\mathcal{D}}
\newcommand{\EE}[0]{\mathbb{E}}
\newcommand{\FF}[0]{\mathbb{F}}
\newcommand{\Ff}[0]{\mathcal{F}}
\newcommand{\Ii}[0]{\mathcal{I}}
\newcommand{\Gg}[0]{\mathcal{G}}
\newcommand{\HH}[0]{\mathbb{H}}
\newcommand{\Hh}[0]{\mathcal{H}}
\newcommand{\Kk}[0]{\mathcal{K}}
\newcommand{\Ll}[0]{\mathcal{L}}
\newcommand{\Mm}[0]{\mathcal{M}}
\newcommand{\NN}[0]{\mathbb{N}}
\newcommand{\Nn}[0]{\mathcal{N}}
\newcommand{\PP}[0]{\mathbb{P}}
\newcommand{\Pp}[0]{\mathcal{P}}
\newcommand{\QQ}[0]{\mathbb{Q}}
\newcommand{\RR}[0]{\mathbb{R}}
\newcommand{\Rr}[0]{\mathcal{R}}
\newcommand{\Ss}[0]{\mathcal{S}}
\newcommand{\TT}[0]{\mathbb{T}}
\newcommand{\Tt}[0]{\mathcal{T}}
\newcommand{\UU}[0]{\mathbb{U}}
\newcommand{\Uu}[0]{\mathcal{U}}
\newcommand{\Vv}[0]{\mathcal{V}}
\newcommand{\Ww}[0]{\mathcal{W}}
\newcommand{\Xx}[0]{\mathcal{X}}
\newcommand{\Yy}[0]{\mathcal{Y}}
\newcommand{\ZZ}[0]{\mathbb{Z}}
\newcommand{\Zz}[0]{\mathcal{Z}}
\newcommand{\acts}[0]{\circlearrowright}
\newcommand{\explain}[2] {
	\begin{flalign*}
		 && \text{#2} && \text{#1}
	\end{flalign*}
}
\newcommand{\maketitleprint}[1][green]{
	\begin{center}
		\includegraphics[width=6cm]{dragon_#1}
	\end{center}
}

% theorem commands
\newtheoremstyle{c_remark}
	{}	% Space above
	{}	% Space below
	{}% Body font
	{}	% Indent amount
	{\bfseries}	% Theorem head font
	{}	% Punctuation after theorem head
	{.5em}	% Space after theorem head
	{\thmname{#1}\thmnumber{ #2}\thmnote{ \normalfont{\text{(#3)}}}}	% head content
\newtheoremstyle{c_definition}
	{3pt}	% Space above
	{3pt}	% Space below
	{}% Body font
	{}	% Indent amount
	{\bfseries}	% Theorem head font
	{}	% Punctuation after theorem head
	{.5em}	% Space after theorem head
	{\thmname{#1}\thmnumber{ #2}\thmnote{ \normalfont{\text{(#3)}}}}	% head content
\newtheoremstyle{c_plain}
	{3pt}	% Space above
	{3pt}	% Space below
	{\itshape}% Body font
	{}	% Indent amount
	{\bfseries}	% Theorem head font
	{}	% Punctuation after theorem head
	{.5em}	% Space after theorem head
	{\thmname{#1}\thmnumber{ #2}\thmnote{ \text{(#3)}}}	% head content

\ifcsname c@english\endcsname
	\theoremstyle{plain}
	\newtheorem{theorem}{Theorem}[section]
	\newtheorem*{theorem*}{Theorem}
	\newtheorem{lemma}[theorem]{Lemma}
	\newtheorem{proposition}[theorem]{Proposition}
	\newtheorem*{proposition*}{Proposition}

	\theoremstyle{definition}
	\newtheorem{definition}[theorem]{Definition}
	\newtheorem*{definition*}{Definition}
	\newtheorem{example}{Example}[section]
	\newtheorem{exercise}{Exercise}[section]

	\theoremstyle{remark}
	\newtheorem*{remark}{Remark}
	\newtheorem*{solution}{Solution}
	\newtheorem{corollary}[theorem]{Corollary}
	\newtheorem{notation}[theorem]{Notation}
\else
	\theoremstyle{c_plain}
	\newtheorem{theorem}{משפט}[section]
	\newtheorem*{theorem*}{משפט}
	\newtheorem{lemma}[theorem]{למה}
	\newtheorem{proposition}[theorem]{טענה}
	\newtheorem*{proposition*}{טענה}

	\theoremstyle{c_definition}
	\newtheorem{definition}[theorem]{הגדרה}
	\newtheorem*{definition*}{הגדרה}
	\newtheorem{example}{דוגמה}[section]
	\newtheorem{exercise}{תרגיל}[section]

	\theoremstyle{c_remark}
	\newtheorem*{remark}{הערה}
	\newtheorem*{solution}{פתרון}
	\newtheorem{corollary}[theorem]{מסקנה}
	\newtheorem{notation}[theorem]{סימון}
\fi

% Questions related commands
\newcounter{question}
\setcounter{question}{1}
\newcounter{sub_question}
\setcounter{sub_question}{1}

\ifcsname c@english\endcsname
	\newcommand{\question}[1][0]{
		\ifthenelse{#1 = 0}{}{\setcounter{question}{#1}}
		\section{Question \arabic{question}}
		\addtocounter{question}{1}
		\setcounter{sub_question}{1}
	}

	\newcommand{\subquestion}[1][0]{
		\ifthenelse{#1 = 0}{}{\setcounter{sub_question}{#1}}
		\subsection{Part \alph{sub_question}}
		\addtocounter{sub_question}{1}
	}
\else
	\newcommand{\question}[1][0]{
		\ifthenelse{#1 = 0}{}{\setcounter{question}{#1}}
		\section{שאלה \arabic{question}}
		\addtocounter{question}{1}
		\setcounter{sub_question}{1}
	}

	\newcommand{\subquestion}[1][0]{
		\ifthenelse{#1 = 0}{}{\setcounter{sub_question}{#1}}
		\subsection{סעיף \localecounter{letters.gershayim}{sub_question}}
		\addtocounter{sub_question}{1}
	}
\fi

% import lua and start of document
\directlua{common = require ('../common')}

\GetEnv{AUTHOR}

% headers
\author{\AUTHOR}
\date\today

\title{פתרון מטלה 04 --- יסודות באלגברה קומוטטיבית}

\DeclareMathOperator{\Ann}{Ann}

\begin{document}
\maketitle
\maketitleprint[red]

\question{}
יהי חוג ויהי $M$ מודול מעליו, נראה ש־$M$ נוצר סופית אם ורק אם הוא איזומורפי למודול מנה של $R^n$ לאיזשהו $0 < n \in \NN$.
\begin{proof}
	נניח ש־$M$ נוצר סופית ונסמן $M = \langle m_1, \ldots, m_n \rangle$.
	נסמן $S \subset R[x_1, \ldots, x_n]$ אוסף הפולינומים כך שמתקיים $f(m_1, \ldots, m_n) = 0 \in M$ לכל $f \in S$.
	נגדיר גם את $I = Z(S)$, אז $I \subset R^n$ אידאל, ונראה ש־$M \cong R^n / I$.
	נגדיר את $\varphi : M \to R^n$ על־ידי $\varphi(m_i) = e_i$ לכל $1 \le i \le n$ ועבור $e_i$ הבסיס הסטנדרטי.
	ברור כי $\im \varphi = R^n$, ונותר להראות ש־$\ker \varphi = I$, אך זוהי למעשה בדיוק הגדרת $I$, לכן ממשפט האיזומורפיזם הראשון נקבל את הדרוש.

	בכיוון ההפוך נניח ש־$M = R^n / I$ ונראה שהוא נוצר סופית.
	נסמן $a_1 = (1, 0, \ldots, 0), \ldots, a_n = (0, \ldots, 0, 1)$, ויהי $b + I \in M$, כאשר $b \in R^n$ נציג, אז ברור כי $b = \sum_{i = 1}^n b_i a_i$, כלומר המודול אכן נוצר סופית.
\end{proof}

\question{}
יהי $R$ חוג ויהי $M$ מודול מעליו, ויהי $N \subset M$ תת־מודול.
נניח ש־$N, M / N$ נוצרים סופית, ונראה שגם $M$ נוצר סופית.
\begin{proof}
	נסמן $N = \langle a_1, \ldots, a_n \rangle$ וכן $M / N = \langle b_1, \ldots, b_m \rangle$.
	יהי $m \in M$, אז $m = \sum \beta_i b_i + n$ עבור $n \in N$, אבל $n = \sum \alpha_i a_i$ לאיזשהו צירוף, אז נקבל בדיוק את המבוקש.
\end{proof}

\question{}
נאמר ש־$R$־מודול $M \ne 0$ הוא פשוט אם ורק אם אין לו תת־מודולים לא טריוויאליים.
נגדיר את המאפס של $M$ להיות,
\[
	\Ann_R M
	= \{ r \in R \mid \forall m \in M, r m = 0 \}
\]
ונאמר ש־$M$ הוא נאמן אם $\Ann_R M = 0$.

\subquestion{}
נראה ש־$\Ann_R M \triangleleft R$ אידאל לכל מודול מעל $R$.
\begin{proof}
	עלינו להראות שזוהי חבורה חיבורית וכן סגירות לכפל ב־$R$.
	ברור כי $0 \in \Ann_R M$, וכן אם $r \in \Ann_R M$ אז לכל $m \in M$ מתקיים $r m = 0$, וכן $0 m = 0$, אז $0 = 0 m = (r - r) m = r m + (- r) m = 0 + (-r) m$ ולכן גם $-r \in \Ann_R M$.
	נסיק שהמאפס הוא חבורה חיבורית, נותר להראות שהוא סגור לכפל.
	יהי $r \in \Ann_R M$, ויהי $s \in R$, אז לכל $m \in M$ מתקיים $(r s) m = s (r m) = 0$ וקיבלנו את המבוקש.
\end{proof}

\subquestion{}
נראה שלכל חוג יש מודול פשוט.
\begin{proof}
	יהי $I \triangleleft R$ אידאל מקסימלי כלשהו, אז $M = R / I$ הוא מודול וכן גם שדה, נראה שהוא מודול פשוט.
	נניח ש־$N \subset M$ תת־מודול. אם לא קיים $0 \ne x \in M$ אז סיימנו, אחרת נניח ש־$x$ הוא ערך כזה, ולכן קיים לו הופכי $x^{-1} \in R / I$,
	כלומר קיים נציג $y \in R$ כך ש־$x \pi(y) = 1$, בפרט $y x \in N$ כמודול ולכן $N = M$.
\end{proof}

\subquestion{}
נניח שקיים $R$־מודול $M$ פשוט ונאמן, נראה ש־$R$ שדה.
\begin{proof}
	נגדיר את ההומומורפיזם $\varphi : R \to M$ על־ידי $\varphi(r) = r \cdot 1_M$.
	יהיו $r, s \in R$ כך ש־$\varphi(r) = \varphi(s)$ אז $(r - s) 1 = 0$, אבל $M$ נאמן ולכן $r - s = 0 \implies r = s$, כלומר $\varphi$ חד־חד ערכי.
	אילו קיים $m \in M$ כך ש־$m \notin \im \varphi$ אז $\varphi(R) \subset M$ הוא תת־מודול לא טריוויאלי, לכן נסיק שלא קיים כזה ו־$\varphi$ הוא על.
	מצאנו ש־$\varphi$ איזומורפיזם, בפרט הוא מעיד כי אין מחלקי אפס ואין אידאלים לא טריוויאליים, ולכן $R$ שדה.
\end{proof}

\question{}
יהי $R$ חוג ויהיו $I, J \triangleleft R$ אידאלים,
נראה כי,
\[
	0
	\to I \cap J
	\xrightarrow{x \mapsto (x, -x)} I \oplus J
	\to I + J
	\to 0
\]
סדרה מדויקת.
\begin{proof}
	התנאי היחיד שמתקיים הוא שלכל $i$ מתקיים, $\im \varphi_i = \ker \varphi_{i + 1}$. \\
	ברור כי $\im \varphi_0 = 0$ ולכן נראה ש־$\ker \varphi_1 = \{ 0 \}$, אך זה ברור מהגדרה והעובדה ש־$\pi_1 \circ \varphi_1$ הוא שיכון $I \cap J \hookrightarrow I$.
	מתקיים $\im \varphi_1 = \{ (x, -x) \mid x \in I \cap J \}$ ולכן נראה שזהו בדיוק $\ker \varphi_2$.
	מתקיים $(x, y) \in \ker \varphi_2 \iff x + y = 0 \implies x = -y$ ולכן שוב קיבלנו שוויון.
	ברור כי $\varphi_2$ היא על, כלומר $\im \varphi_2 = I + J$, וברור כי $\ker \varphi_3 = I + J$.
\end{proof}

\question{}
יהיו $M, M', M''$ מודולים מעל $R$, ויהיו $v : M \to M'', u : M' \to M$ הומומורפיזמים. \\
לכל $R$־מודול $N$ נגדיר את $\bar{u} : \Hom_R(N, M') \to \Hom_R(N, M)$ על־ידי, $\bar{u}(\varphi) = u \circ \varphi$ וכן נגדיר את $\bar{v} : \Hom_R(N, M) \to \Hom_R(N, M'')$ על־ידי $\bar{v}(\varphi) = v \circ \varphi$. \\
נראה שהסדרה,
\[
	0
	\to M'
	\xrightarrow{u} M
	\xrightarrow{v} M''
\]
מדויקת אם ורק אם,
\[
	0
	\to \Hom_R(N, M')
	\xrightarrow{\bar{u}} \Hom_R(N, M)
	\xrightarrow{\bar{v}} \Hom_R(N, M'')
\]
מדויקת.
\begin{proof}
	שקול לתרגיל שהוצג בתרגול.
\end{proof}

\question{}
בשאלה זו נוכיח שאם $A$ תחום ראשי ו־$M$ מודול חופשי, אז כל $N \subset M$ הוא חופשי. \\
יהי ${\{ e_i \}}_{i \in I}$ בסיס ל־$M$, ויהי $\le$ סדר טוב על $I$, וכן נסמן $p_i : M \to A$ ההטלה לקורדינטה $i$. \\
לכל $i$ נסמן $M_i = (e_j \mid j \le i)$ וכן נסמן $N_i = N \cap M_i$.

\subquestion{}
נראה ש־$N = \bigcup_{i \in I} N_i$.
\begin{proof}
	מספיק להראות ש־$M = \bigcup M_i$, ונקבל שהטענה נובעת.
	יהי $x \in M$, אז $p_i(x) = 0$ לכמעט כל $i$, ונסמן את $n = \max\{ i \in I \mid p_i(x) \ne 0 \}$, נבחין כי הגדרה זו תקפה מהעובדה שזהו מקסימום של קבוצה סופית.
	בהתאם נובע ש־$x \in \bigcup_{i \le n} M_i \subseteq \bigcup_{i \in I} M_i$.
\end{proof}

\subquestion{}
נראה ש־$p_i(N_i) \triangleleft A$ אידאל של $A$.
\begin{proof}
	ברור כי $0 \in p_i(N_i)$, וכן אם $x \in N_i$ אז גם $-1 \cdot p_i(x) = p_i(-x) \in N_i$ ולכן זוהי חבורה חיבורית.
	יהי $a \in A$, ונניח ש־$x = \sum_{i \in I} \alpha_i x_i$, אז,
	\[
		a x
		= a \sum_{i \in I} \alpha_i x_i
		= \sum_{i \in I} (a \alpha_i) x_i
		\implies p_i(a x)
		= a p_i(x)
	\]
	ולכן זהו אכן אידאל.
\end{proof}

לכל $i \in I$ נבחר $a_i \in A$ כך ש־$p_i(N_i) = (a_i)$, וכן נבחר $n_i \in N_i$ כך ש־$p_i(n_i) = a_i$, בפרט אם $a_i = 0$ אז נבחר $n_i = 0$.
\subquestion{}
נראה שמתקיים $(n_i \mid i \in I) = N$.
\begin{proof}
	מסעיף א' מספיק להראות שלכל $i \in I$ מתקיים $N_i \subset (n_j \mid j \in I)$, לכן נניח בשלילה שזה לא מתקיים, ויהי $i$ המינימלי עבורו $N_i \not\subset (n_j \mid j \in I)$.
	נסמן $x \in N_i \setminus (n_j \mid j \in I)$, בפרט $x_i = p_i(x) \in (a_i)$ ולכן $x_i = \alpha a_i$, ונסמן $y = x - \alpha n_i$, לכן $y \in N_i$ וגם $x \in (n_j \mid j \in I) \iff y \in (n_j \mid j \in I)$.
	לבסוף נובע ש־$y \notin (n_j \mid j \in I)$ אבל $y \in N_{i - 1}$, כלומר $N_{i - 1} \not\subset (n_j \mid j \in I)$ בסתירה למינימליות $i$.
\end{proof}

\subquestion{}
נראה ש־$\{ n_i \mid i \in I, n_i \ne 0 \}$ היא בלתי־תלויה לינארית.
\begin{proof}
	נניח ש־$n_i \in (n_j \mid j \in I \setminus \{ i \})$, ועל־ידי החלפת אגפים נוכל להניח ש־$i$ מקסימלי (שוב מטענת סופיות המקדמים השונים מאפס).
	נסמן $i' < i$ כאינדקס הגדול ביותר כך ש־$n_{i'} \ne 0$, לכן $N_i = N_{i'}$, כלומר $n_i = \alpha n_{i'}$ לאיזשהו $\alpha \in A$.
	בפרט נסיק ש־$p_i(n_i) = p_i(\alpha n_{i'})$ בסתירה.
\end{proof}

\subquestion{}
נסיק ש־$N$ מודול חופשי.
\begin{proof}
	מצאנו ש־$n_i$ פורשים את $N$ וכן שאין תלות, כלומר יש הצגה יחידה, לכן $N$ מודול חופשי.
\end{proof}

\question{}
נראה ש־$1 \otimes i + i \otimes 1 \in \CC \otimes_{\RR} \CC$ לא טנזור פשוט.
\begin{proof}
	נסמן $z = a + i b, w = c + i d$ ולכן ונניח ש־$z \otimes w = 1 \otimes i + i \otimes 1$, אבל,
	\[
		z \otimes w
		= (a + i b) \otimes (c + i d)
		= (a \otimes c) + (i b \otimes i d) + (a \otimes i d) + (i b \otimes c)
		= a c (1 \otimes 1) + b d (i \otimes 1) + a d (1 \otimes i) + b c (i \otimes 1)
	\]
	ולכן נקבל שבהכרח $a c = 0, b d = 0, a d = 1, b c = 1$.
	נובע ש־$a, b, c, d \ne 0$ אבל $a c = 0$, כלומר קיבלנו סתירה לטענה שזהו טנזור פשוט.
\end{proof}

\question{}
יהיו $N, M$ מודולים מעל $R$, ויהי $I \triangleleft R$ אידאל.
נראה שמתקיים,
\[
	(M / I M) \otimes_R (N / I N)
	\cong (M \otimes_R N) / (I (M \otimes_R N))
\]
\begin{proof}
	נגדיר את $\varphi : (M / I M) \times (N / I N) \to (M \otimes_R N) / (I (M \otimes_R N))$ על־ידי,
	\[
		\varphi(m + I M, n + I N)
		= (m \otimes n) + I (M \otimes_R N)
	\]
	כאשר ביטוי זה מוגדר היטב מבילינאריות המכפלה הטנזורית.
	לכן נובע שקיים הומומורפיזם,
	\[
		\tilde{\varphi} : (M / I M) \otimes_R (N / I N) \to (M \otimes_R N) / (I (M \otimes_R N))
	\]
	מהצד השני נגדיר את,
	\[
		\psi : (M \otimes_R N) / (I (M \otimes_R N)) \to (M / I M) \otimes_R (N / I N),
		\qquad
		\psi((m \otimes n) + I (M \otimes_R N))
		= (m + I M) \otimes (n + I N)
	\]
	שוב ההעתקה מוגדרת היטב על־ידי שימוש בבילינאריות.
	ברור כי ההומומורפיזמים הפוכים ולכן נסיק שהם איזומורפיזמים, ובפרט הזהות נכונה.
\end{proof}

\question{}
יהי $M$ מודול מעל $R$ ונסמן $M[x]$ את הפולינומים עם מקדמים ב־$M$.
נגדיר את $M[x]$ כמודול מעל $R$ על־ידי הפעולה,
\[
	(m_0 x^0 + \cdots + m_n x^n) + (m_0' x^0 + \cdots + m_k' x^k)
	= \sum_{i = 0}^{\max\{ n, m \}} (m_i + m_i') x^i
\]
והפעולה,
\[
	a (m_0 x^0 + \cdots + m_n x^n)
	= (a m_0) x^0 + \cdots + (a m_n) x^n
\]
נראה כי מתקיים,
\[
	M \otimes_R R[x]
	\cong M[x]
\]
\begin{proof}
	נגדיר את ההעתקה,
	\[
		\varphi : M \times R[x] \to M[x],
		\qquad
		\varphi(m, r_0 x^0 + \cdots + r_n x^n)
		= (r_0 m) x^0 + \cdots + (r_n m) x^n
	\]
	נבחין מבדיקה ישירה כי זהו הומומורפיזם מודולים, ולכן יש לו הרמה להומומורפיזם,
	\[
		\tilde{\varphi} : M \otimes_R R[x] \to M[x]
	\]
	מהצד השני נגדיר את,
	\[
		\psi : M[x] \to M \otimes_R R[x],
		\qquad
		\psi(\sum_{i = 0}^n m_i x^i)
		= \sum_{i = 0}^n m_i \otimes x^i
	\]
	ולכן גם $\psi$ היא פונקציה מוגדרת היטב ומבדיקה ישירה היא גם הומומורפיזם, ונרצה להראות שהיא הפוכה ל־$\tilde{\varphi}$,
	\begin{align*}
		\psi(\tilde{\varphi}(m \otimes (r_0 x^0 + \cdots + r_n x^n)))
		& = \psi(\varphi(m, r_0 x^0 + \cdots + r_n x^n)) \\
		& = \psi((r_0 m) x^0 + \cdots + (r_n m) x^n) \\
		& = ((r_0 m) \otimes x^0) + \cdots + ((r_n m) \otimes x^n) \\
		& = m \otimes (r_0 x^0 + \cdots + r_n x^n)
	\end{align*}
	ומהצד השני,
	\[
		\tilde{\varphi}(\psi(m_0 x^0 + \cdots + m_n x^n))
		= \tilde{\varphi}(m_0 \otimes x^0 + \cdots + m_n \otimes x^n)
		= \tilde{\varphi}(m_0 \otimes x^0) + \cdots + \tilde{\varphi}(m_n \otimes x^n)
		= m_0 x^0 + \cdots + m_n x^n
	\]
	וקיבלנו בדיוק שהפונקציות הופכיות ולכן איזומורפיזמים ומעידות על הרצוי.
\end{proof}

\question{}
יהיו $M, N, P$ מודולים מעל $R$, נראה כי,
\[
	\Hom_R(M \otimes_R N, P)
	\cong \Hom_R(M, \Hom_R(N, P))
\]
\begin{proof}
	ברור כי הטענה נכונה עבור מכפלות סטנדרטיות, ושאר הטענה נובעת מקיבוע הומומורפיזם ושימוש בתכונת ההרמה.
\end{proof}

\end{document}
